% !TeX program = lualatex
% !TeX encoding = utf8
% !TeX spellcheck = uk_UA
% !TeX root = ../EMConspect.tex
% УВАГА: для \enquote{...} потрібен пакет csquotes (\usepackage[ukrainian=quotes]{csquotes}
% в преамбулі EMConspect.tex) --- переконайтесь, що він підключений.

%=========================================================
\chapter{Мікроскопічна теорія діелектричної проникності}\label{\currfilebase-II}
\graphicspath{{\currfilebase/Pictures}}
%=========================================================
\epigraph{Макроскопічна діелектрична проникність --- це лише тінь, яку відкидає на вимірювальний прилад складна динаміка мікроскопічних диполів.
}{Переосмислення ідей сучасної фізики}

У попередньому розділі ми ввели діелектричну проникність $\epsilon$ як \emphx{феноменологічну} величину: ми знаємо, як її виміряти (через ємність
конденсатора), і як нею користуватись у розрахунках, але поки що не знаємо, \emphx{звідки} вона береться.

%% --------------------------------------------------------
\section{Механізми поляризовності молекули}\label{sec:AtomPolarizability}
%% --------------------------------------------------------

У зовнішньому полі $\Efield$ молекула набуває наведеного дипольного моменту $\vect p = \alpha\Efield$ (формула~\eqref{eq:InducedDipole}). Але
\emphx{звідки} береться цей момент --- питання не одне, а фізично різне залежно від будови молекули:

\begin{itemize}
	\item $\alpha_e$ --- \emphx{електронна}: поле зміщує легку електронну хмару відносно ядра. Є \emphx{в усіх} атомах і молекулах без винятку,
	      встигає за полем аж до оптичних частот ($\sim10^{15}$ Гц);
	\item $\alpha_i$ --- \emphx{йонна}: у йонних кристалах поле зміщує одні йони відносно інших (маса набагато більша за електронну) ---
	      \enquote{вимикається} вже в ІЧ-діапазоні;
	\item $\alpha_\text{ор}$ --- \emphx{орієнтаційна}: якщо молекула \emphx{вже має} власний дипольний момент $p_0$ (полярна молекула), поле не
	      наводить, а \emphx{розвертає} його. Найповільніший і найчутливіший до $T$ механізм --- \enquote{вимикається} вже в НВЧ-діапазоні.
\end{itemize}

Повна статична поляризовність --- сума всіх діючих на цій частоті внесків, $\alpha = \alpha_e + \alpha_i + \alpha_\text{ор}$.


%% --------------------------------------------------------
\section{Орієнтаційна поляризація: закон Ланжевена--Дебая}\label{sec:Langevin-Debye}
%% --------------------------------------------------------

Полярна молекула ($p_0\neq0$) в полі $\Efield$ має потенціальну енергію
\begin{equation*}
	U(\theta) = -\vect p_0\cdot\Efield = -p_0 E\cos\theta
\end{equation*}
(п.~\ref{sec:DipoleEnergy}), де $\theta$ --- кут між диполем і полем; енергія мінімальна ($U=-p_0E$) при орієнтації вздовж поля, максимальна
($U=+p_0E$) --- проти поля. За відсутності поля тепловий рух орієнтує диполі цілком хаотично; у полі ж число молекул, орієнтованих в одиничний
тілесний кут навколо напрямку $\theta$, підкоряється \emphx{розподілу Больцмана}:
\begin{equation}\label{eq:Boltzmann}
	n(\theta) = n_0\, e^{-U/kT} = n_0\, e^{+p_0 E\cos\theta / kT}.
\end{equation}
Це й означає порушення хаосу: молекул, орієнтованих вздовж поля ($\cos\theta=1$, менша енергія), \emphx{експоненційно більше}, ніж орієнтованих
проти поля ($\cos\theta=-1$).

У практично важливому випадку --- не надто сильне поле й не надто низька температура, $p_0E \ll kT$ --- показник експоненти малий, тож
розподіл~\eqref{eq:Boltzmann} можна розкласти в ряд і обмежитись лінійним членом:
\begin{equation}\label{eq:BoltzmannLinear}
	n(\theta) \approx n_0\left(1 + \frac{p_0 E \cos\theta}{kT}\right).
\end{equation}
Фізично: за кімнатної температури $kT$ на багато порядків перевищує енергію орієнтації $p_0E$ навіть у сильних лабораторних полях, тому переважна
орієнтація --- \emphx{слабка} поправка до майже хаотичного розподілу, а не повне вишикування диполів вздовж поля.

Поляризація $P$ --- це дипольний момент одиниці об'єму. У напрямку, перпендикулярному полю, внески різних диполів у середньому компенсуються (за
симетрією $\theta\to\pi-\theta$); вздовж поля --- ні. Тому $P$ напрямлена вздовж поля і дорівнює добутку концентрації молекул $n$ на середню
проекцію дипольного моменту:
\begin{equation*}
	P = n\, p_0\, \overline{\cos\theta} = n p_0 \cdot \frac{\displaystyle\int_0^\pi \cos\theta \, n(\theta)\, 2\pi\sin\theta\, d\theta}
	{\displaystyle\int_0^\pi n(\theta)\, 2\pi\sin\theta\, d\theta}.
\end{equation*}
Підставляючи сюди лінеаризований розподіл~\eqref{eq:BoltzmannLinear} і виконуючи інтегрування за кутом (у знаменнику лишається $n_0$, у
чисельнику --- лише доданок з $\cos\theta$, оскільки $\int\cos\theta\,d\Omega=0$, а $\int\cos^2\theta\,d\Omega = 4\pi/3$), одержуємо
\begin{equation}\label{eq:LangevinDebye}
	\tcbhighmath{
		P = \frac{n p_0^2}{3kT}\, E.
	}
\end{equation}
Це і є \emphx{закон Ланжевена--Дебая}: поляризація речовини лінійна за полем, а коефіцієнт пропорційності явно виражений через мікроскопічні
характеристики молекули ($p_0$) і речовини ($n$, $T$).

Оскільки за означенням $P = \chi E$, порівняння з~\eqref{eq:LangevinDebye} одразу дає поляризовність речовини (а заразом і орієнтаційну
поляризовність \emphx{однієї} молекули, $\chi = n\alpha_\text{ор}$):
\begin{equation*}
	\chi = \frac{n p_0^2}{3kT}
	\qquad\Longrightarrow\qquad
	\alpha_\text{ор} = \frac{p_0^2}{3kT}.
\end{equation*}
Підставляючи $\chi$ у означення проникності $\epsilon = 1+4\pi\chi$, знаходимо остаточний макроскопічний результат:
\begin{equation}\label{eq:CurieLaw}
	\tcbhighmath{
		\epsilon = 1 + \frac{4\pi n p_0^2}{3kT}.
	}
\end{equation}

\begin{Attention}
	$\alpha_\text{ор} \sim 1/T$ --- \emphx{закон Кюрі}: чим вища температура, тим інтенсивніший тепловий рух, який руйнує впорядкування диполів.
	Якщо $\epsilon$ речовини помітно спадає з ростом $T$ --- пряма ознака, що діелектрик полярний (електронна й йонна поляризовність від $T$
	практично не залежать, п.~\ref{sec:AtomPolarizability}).
\end{Attention}

%% --------------------------------------------------------
\section{Локальне поле. Формула Клаузіуса--Моссотті}\label{sec:ClausiusMossotti}
%% --------------------------------------------------------

Для розріджених середовищ справджується наївне співвідношення $\chi = n\alpha$: кожна молекула відчуває лише зовнішнє (середнє) поле $\Efield$. У
густому ж середовищі кожна молекула занурена в поле \emphx{вже поляризованих} сусідів, тож реальне (локальне) поле $\Efield_\text{лок}$, що на неї
діє, --- не $\Efield$, а деяка відмінна від нього величина, яку й належить знайти.

Ми знаємо, що поле всередині однорідно поляризованої кулі однорідне й дорівнює $-\frac{4\pi}{3}\vect P$ (деполяризаційне
поле). Оточимо обрану молекулу уявною сферою Лоренца (мала порівняно з тілом, велика порівняно з міжмолекулярними відстанями). Поле на молекулу ---
сума трьох внесків (рис.~\ref{pic:CavityField}): середнє поле $\Efield$; мінус поле кулі-вирізу, $+\frac{4\pi}{3}\vect P$; і поле дискретних
сусідів усередині сфери. Для кубічних кристалів, газів і слабополярних рідин останній внесок за симетрією дорівнює нулю, звідки

\begin{figure}[h!]\centering
	\begin{tikzpicture}[>=latex, scale=0.8]
		\draw[thick, name path=blob]
			plot [smooth cycle] coordinates {
				(-2.2,-1.8) (-1.0,-2.3) (0.4,-2.1) (1.6,-2.4)
				(2.4,-1.4) (2.6,-0.1) (2.0,1.1) (2.5,2.0)
				(1.2,2.5) (0.0,2.0) (-1.3,2.4) (-2.4,1.2)
				(-2.7,0.0) (-2.3,-1.0)
			};
		\begin{scope}
			\clip plot [smooth cycle] coordinates {
				(-2.2,-1.8) (-1.0,-2.3) (0.4,-2.1) (1.6,-2.4)
				(2.4,-1.4) (2.6,-0.1) (2.0,1.1) (2.5,2.0)
				(1.2,2.5) (0.0,2.0) (-1.3,2.4) (-2.4,1.2)
				(-2.7,0.0) (-2.3,-1.0)
			};
			\foreach \d in {-4.5,-4.2,...,4.5} {
				\draw[gray!70, thin] (\d,-3) -- ++(1,1) -- ++(2.5,2.5);
			}
			\fill[white] (0,0) circle (1.0);
		\end{scope}
		\draw[thick, dashed] (0,0) circle (1.0);
		\foreach \a in {55,70,90,110,125} {
			\node[font=\scriptsize] at ({1.15*cos(\a)},{1.15*sin(\a)}) {$-$};
		}
		\foreach \a in {235,250,270,290,305} {
			\node[font=\scriptsize] at ({1.15*cos(\a)},{1.15*sin(\a)}) {$+$};
		}
		\draw[->, thick, brown] (0,1.35) -- ++(0,0.7) node[right] {$\vect P$};
		\draw[->, thick, blue] (0,-0.4) -- ++(0,0.75) node[right] {$\Efield_\text{лок}$};
		\node[circle, fill=green!70!black, inner sep=0, minimum size=0.2cm] at (0,0) {};
	\end{tikzpicture}
	\caption{Сферична порожнина Лоренца в поляризованому діелектрику\label{pic:CavityField}}
\end{figure}

\begin{equation}\label{eq:LorentzField}
	\tcbhighmath{
		\Efield_\text{лок} = \Efield + \frac{4\pi}{3}\vect P.
	}
\end{equation}

\begin{Attention}
	Наближення: працює для кубічних кристалів і більшості неполярних/слабополярних речовин. Для сильно полярних рідин з водневими зв'язками
	(вода) --- потребує уточнення (теорія Онзагера).
\end{Attention}

Маючи локальне поле~\eqref{eq:LorentzField}, запишемо наведений момент молекули саме через нього: $\vect p = \alpha\Efield_\text{лок}$. Тоді для
поляризації (дипольного моменту одиниці об'єму з концентрацією молекул $n$) маємо
\begin{equation*}
	\vect P = n\vect p = n\alpha\Efield_\text{лок} = n\alpha\left(\Efield + \frac{4\pi}{3}\vect P\right).
\end{equation*}
Це --- рівняння відносно самої $\vect P$: вона входить і в ліву, і в праву частину, бо локальне поле саме поляризацією середовища й визначається
(задача самоузгоджена --- поляризація створює поле, яке, своєю чергою, впливає на поляризацію). Розкриваючи дужки та групуючи доданки з $\vect P$:
\begin{equation*}
	\vect P\left(1 - \frac{4\pi n\alpha}{3}\right) = n\alpha\,\Efield
	\quad\Longrightarrow\quad
	\vect P = \frac{n\alpha}{1 - \dfrac{4\pi n\alpha}{3}}\,\Efield.
\end{equation*}
Порівнюючи це з означенням поляризовності речовини $\vect P = \chi\Efield$, знаходимо мікроскопічний вираз для $\chi$, що вже враховує локальне
поле:
\begin{equation}\label{eq:ChiFromAlpha}
	\chi = \frac{n\alpha}{1 - \dfrac{4\pi n\alpha}{3}}.
\end{equation}

Залишилось підставити цей вираз у означення проникності $\epsilon = 1+4\pi\chi$. Зручно скористатись тим, що $\epsilon - 1 = 4\pi\chi$ напряму, а
отже $\epsilon + 2 = (\epsilon-1) + 3 = 4\pi\chi + 3$. Обчислимо чисельник і знаменник майбутнього відношення окремо, звівши кожен до спільного
знаменника $1-\frac{4\pi n\alpha}{3}$:
\begin{multline*}
	\epsilon - 1 = 4\pi\chi = \frac{4\pi n\alpha}{1-\dfrac{4\pi n\alpha}{3}},
	\\
	\epsilon + 2 = 4\pi\chi + 3 = \frac{4\pi n\alpha + 3\left(1-\dfrac{4\pi n\alpha}{3}\right)}{1-\dfrac{4\pi n\alpha}{3}}
	= \frac{4\pi n\alpha + 3 - 4\pi n\alpha}{1-\dfrac{4\pi n\alpha}{3}} = \frac{3}{1-\dfrac{4\pi n\alpha}{3}}.
\end{multline*}
Однакові знаменники в чисельнику й знаменнику відношення $(\epsilon-1)/(\epsilon+2)$ скорочуються, і залишається напрочуд компактний та симетричний
результат:
\begin{equation}\label{eq:ClausiusMossotti}
	\tcbhighmath{
		\frac{\epsilon - 1}{\epsilon + 2} = \frac{4\pi}{3}\, n\alpha.
	}
\end{equation}
Це --- \emphx{формула Клаузіуса--Моссотті}. Її головна цінність --- вона \emphx{напряму} зв'язує макроскопічно вимірювану величину ($\epsilon$) з
мікроскопічною молекулярною характеристикою ($\alpha$) і концентрацією $n$, причому враховує (на відміну від наївної $\chi = n\alpha$) ефекти
взаємної поляризації сусідніх молекул. Використовують її в обидва боки: за відомим $\alpha$ обчислюють $\epsilon$, або, навпаки, за виміряним
$\epsilon$ й відомою концентрацією $n$ знаходять $\alpha$ --- один з небагатьох прямих макроскопічних методів вимірювання мікроскопічної
характеристики речовини.

Для розрідженого газу ($\frac{4\pi}{3}n\alpha\ll1$) знаменник у~\eqref{eq:ChiFromAlpha} можна замінити одиницею, і формула зводиться до наївного
$\chi\approx n\alpha$, з якого ми й почали; поправка на локальне поле стає суттєвою вже в густій рідині чи твердому тілі.

\begin{Attention}
	\emphx{Оцінка.} $\alpha \sim 10^{-24}\ \text{см}^3$ (атомний об'єм). Газ ($n\sim2{,}7\times10^{19}\ \text{см}^{-3}$): поправка $\sim10^{-4}$ ---
	нехтовна. Рідина ($n\sim3\times10^{22}\ \text{см}^{-3}$): поправка $\sim10\%$ --- нехтувати не можна.
\end{Attention}

Наостанок --- місток до оптики: на оптичних частотах поляризовність лише електронна (\S\ref{sec:AtomPolarizability}), а $\epsilon_\text{опт} =
n_\text{зал}^2$, тож
\begin{equation}\label{eq:LorentzLorenz}
	\frac{n_\text{зал}^2 - 1}{n_\text{зал}^2 + 2} = \frac{4\pi}{3}\, n\alpha_e
\end{equation}
--- \emphx{формула Лоренца--Лоренца}, основа поняття молярної рефракції.